% !PRINTED 410
\weberpage{410}
%<<CONT
functions can be constant, then by (6) Theorem 4.
is proved.

We apply this theorem to the function $\vartheta$ occurring in 3.
If we perform with the variables $x$ any permutation
of the alternating group, the functions
$u, u_1, \ldots, u_{\nu-1}$ likewise undergo a permutation. By
(5), $\vartheta$ passes under this permutation into another function $\vartheta_1$,
which by (4) is rationally expressible in terms of $\vartheta$. But
$\vartheta$ is likewise rationally expressible in terms of $\vartheta_1$, since in
Theorem 3. one may replace $\vartheta$ by $\vartheta_1$, and $\vartheta$ also arises from $\vartheta_1$
by a permutation of the alternating group. Hence
by 4., $\vartheta$ and $\vartheta_1$ are \gesperrt{linear functions} of one another.
From this the following result is obtained:

5.~\gesperrt{If for a general equation of degree $n$
a resolvent with one parameter exists,
then there is a rational function $\vartheta$
of $n$ variables $x$ which, under the permutations
of the alternating group, passes into a linear function
of itself}
\weq{7}{\chi(\vartheta) = \frac{a\vartheta + b}{c\vartheta + d}}
\gesperrt{where $a, b, c, d$ are constants, i.\,e.\ numbers.}

\begin{center}
\S.\,109.

\gesperrt{Group of the resolvents with one parameter.}
\end{center}

The function $\vartheta$ that occurs in Theorem 5. of the preceding section
is, by Theorem 3., itself the root of a resolvent
with \gesperrt{one parameter} $\Phi(\vartheta, z) = 0$. This equation,
since every root is rationally expressible in terms of every other,
is a normal equation, and its \gesperrt{Galois} group is isomorphic
with the group of the linear substitutions $\chi(\vartheta)$ of
Theorem 5. Since one can now carry out in the identical equation $\Phi(\vartheta, z) = 0$
all the permutations of the alternating group $A$ of the variables $x$,
whereby $z$ does not change, while $\vartheta$ passes into every
other root of $\Phi$, the group of the linear
substitutions $\chi(\vartheta)$, i.\,e.\ the \gesperrt{Galois} group of the equa-
%CONT>>
